% Drawing 154/04-12 -- folder 154, batch 04, archive page 72 (PDF page 73).
% Transcription: croquis au crayon, la feuille tenue à l'envers du listing : un polygone régulier (octogone) rayonné depuis son centre, une ligne brisée en zigzag épaissie le long d'une partie de son bord ; des faisceaux de droites concourantes, isolés ou sur une droite, avec les légendes « $n - \nu$ » et « $(n-1) - (\nu-1) - (\nu'-1)$ » ; un angle portant deux points ; une pseudo-droite fusiforme esquissée au-dessous
% Kind: pencil sketches: rayed octagon with zigzag, pencils of lines. Method 2 (schematic redraw from the facsimile). Status: partial.
% Notes: Faint pencil page. Drawn: the octagon with its eight rays and centre point, the thick zigzag along its lower left boundary; a star of three lines on its right; a line with a marked point and a transversal at the upper right; a star of four lines on the left; a short line with a marked point; the line carrying two pencils of four lines, with the two legends above it; an angle with two points on its lower side; the two faint curves of the spindle below. The small pencils of lines drawn at the corners of the zigzag are left out; positions are read off a faint page and are approximate. The two legends are read « n - nu » and « (n-1) - (nu-1) - (nu'-1) »: the transcription's note has « n - 1 » and capital N, N' (marked there as partly read); the same two quantities are written with nu, nu' on the drawings of pp. 83 and 93 (« n - nu », « n - 1 - (nu - 1) - (nu' - 1) »), which settles the reading. Decided by a machine pass on the facsimile, 2026-11-12; not checked by a reader. Redrawn 2026-11-12 (issue 128): to check again. 2026-11-14, Claude Opus 5.5, overlay check (issue 147): rasterised from its own coordinates over the 400 dpi colour scan on a unit grid and measured against the colour masks of the ink; not compiled (no TeX in the session); the redraw is unchanged. Upside-down crop. The octagon, the stars and the line with its two pencils lie near their ink, but the faint strokes of the page are only approximately placed: the zigzag along the octagon's lower left boundary and the small pencils at its corners, the angle with two points, and the spindle below (drawn as a symmetric lens, where the page has two faint, uneven curves and small pencil fans) lie up to 3 to 5 units off or are left out. Stays partial: the page needs redrawing stroke by stroke.
% Coordinates: 1 unit = 1 mm = 1/100 of the crop width (bbox in page fractions: [0.21, 0.085, 0.71, 0.435]).
% Redrawn by Claude (claude-fable-5-2), issue #128, 2026-11-12; not checked by a human; not compiled (no TeX in the session).
% Checked by Claude (claude-opus-5-5) 2026-11-14 (issue #147): overlaid on the 400 dpi colour facsimile; not compiled (no TeX in the session). Not checked by a human. Status: partial.
% Rights: reproduces part of the work -- behind RENDER_TRANSCRIPTS, never published by itself.
\documentclass[tikz,border=3pt]{standalone}
\usepackage{amsmath,amssymb}
\usetikzlibrary{patterns,arrows.meta}
\tikzset{pencil/.style={black!45}, struck/.style={black!35}, mpt/.style={circle,draw,inner sep=1.3pt}, pt/.style={circle,fill,inner sep=0.9pt}, opt/.style={circle,draw,fill=white,inner sep=1.1pt}}
\begin{document}
\begin{tikzpicture}[x=1mm,y=1mm,line cap=round,line join=round,font=\scriptsize]
\draw[pencil] (20,-10.5) -- (23,-5.5) -- (32,-5) -- (37.5,-9) -- (38,-15) -- (35,-20) -- (29,-22.5) -- (21,-19) -- cycle;
\draw[pencil] (28,-11.5) -- (20,-10.5);
\draw[pencil] (28,-11.5) -- (23,-5.5);
\draw[pencil] (28,-11.5) -- (32,-5);
\draw[pencil] (28,-11.5) -- (37.5,-9);
\draw[pencil] (28,-11.5) -- (38,-15);
\draw[pencil] (28,-11.5) -- (35,-20);
\draw[pencil] (28,-11.5) -- (29,-22.5);
\draw[pencil] (28,-11.5) -- (21,-19);
\node[pt] at (28,-11.5) {};
\draw[pencil,very thick] (18.5,-18) -- (21,-16) -- (24,-19.5) -- (26,-18) -- (29,-22.5) -- (35,-16.5) -- (37,-17.5);
\draw[pencil] (43,-9) -- (59,-24);
\draw[pencil] (50,-10) -- (52,-23);
\draw[pencil] (44,-18.5) -- (58,-14.5);
\draw[pencil] (63,-13) -- (80,-10.5);
\draw[pencil] (69,-5) -- (75,-20);
\node[pt] at (72.5,-11.6) {};
\draw[pencil] (3,-24) -- (10,-31);
\draw[pencil] (6,-21) -- (7,-34);
\draw[pencil] (1,-29) -- (12,-26);
\draw[pencil] (11,-24) -- (2,-31);
\draw[pencil] (43,-33) -- (62,-32);
\node[pt] at (52,-32.5) {};
\draw[pencil] (68,-30) -- (94,-28.5);
\draw[pencil] (71,-25) -- (76,-34.4);
\draw[pencil] (73,-24.5) -- (74,-34.9);
\draw[pencil] (75.5,-25) -- (71.5,-34.4);
\draw[pencil] (77,-26.5) -- (70,-32.9);
\draw[pencil] (81.5,-25) -- (86.5,-33.2);
\draw[pencil] (83.5,-24) -- (84.5,-34.2);
\draw[pencil] (86,-24.5) -- (82,-33.7);
\draw[pencil] (87.5,-26) -- (80.5,-32.2);
\node[pt] at (73.5,-29.7) {};
\node[pt] at (84,-29.1) {};
\node[inner sep=1pt] at (63,-27) {\small $n-\nu$};
\node[inner sep=1pt] at (72,-22.5) {\small $(n-1)-(\nu-1)-(\nu'-1)$};
\draw[pencil] (39,-30) -- (39.5,-38.5) -- (55,-40);
\node[pt] at (46,-39.1) {};
\node[pt] at (48,-39.3) {};
\draw[pencil] plot[smooth, tension=0.55] coordinates {(52,-44) (48,-57) (47,-70) (50,-85) (53,-95)};
\draw[pencil] plot[smooth, tension=0.55] coordinates {(52,-44) (56,-56) (57,-68) (55,-84) (53,-95)};
\end{tikzpicture}
\end{document}
